\documentclass[11pt,a4paper]{article}

% ── Codificación y fuentes ────────────────────────────────────────────────────
\usepackage[utf8]{inputenc}
\usepackage[T1]{fontenc}
\usepackage{lmodern}
\usepackage[spanish]{babel}

% ── Geometría y espaciado ─────────────────────────────────────────────────────
\usepackage[top=2.5cm, bottom=2.5cm, left=2.8cm, right=2.8cm]{geometry}
\usepackage{setspace}
\setstretch{1.15}
\usepackage{parskip}

% ── Tablas ────────────────────────────────────────────────────────────────────
\usepackage{booktabs}
\usepackage{tabularx}
\usepackage{array}
\usepackage{enumitem}

% ── Colores y cajas ───────────────────────────────────────────────────────────
\usepackage[dvipsnames]{xcolor}
\usepackage{tcolorbox}
\tcbuselibrary{skins, breakable}

% ── Cabeceras y pies de página ────────────────────────────────────────────────
\usepackage{fancyhdr}
\usepackage{lastpage}
\pagestyle{fancy}
\fancyhf{}
\fancyhead[L]{\small\textbf{Informe de Análisis} --- Imágenes}
\fancyhead[R]{\small 2026-06-28}
\fancyfoot[C]{\small Página \thepage\ de \pageref{LastPage}}
\renewcommand{\headrulewidth}{0.4pt}

% ── Hiperenlaces y URLs ──────────────────────────────────────────────────────
\usepackage[colorlinks=true, linkcolor=NavyBlue, urlcolor=NavyBlue]{hyperref}
\usepackage{url}
\def\path#1{\url{#1}}

% ── Permitir saltos de línea más flexibles ───────────────────────────────────
\sloppy
\emergencystretch 3em

% ── Colores personalizados ────────────────────────────────────────────────────
\definecolor{azulTitulo}{HTML}{1B2A6B}
\definecolor{grisFondo}{HTML}{F5F5F5}

% ── Estilos de cajas ─────────────────────────────────────────────────────────
\tcbset{
  cajaTitulo/.style={
    enhanced, breakable,
    colback=azulTitulo!8, colframe=azulTitulo,
    fonttitle=\bfseries\large, coltitle=white,
    attach boxed title to top left={yshift=-2mm, xshift=4mm},
    boxed title style={colback=azulTitulo},
    top=6pt, bottom=6pt, left=8pt, right=8pt,
  },
  cajaContenido/.style={
    enhanced, breakable,
    colback=grisFondo, colframe=gray!50,
    fonttitle=\bfseries, coltitle=black,
    top=4pt, bottom=4pt, left=8pt, right=8pt,
  },
}

% ─────────────────────────────────────────────────────────────────────────────
\begin{document}

% ══ PORTADA ══════════════════════════════════════════════════════════════════
\begin{center}
  {\Large\bfseries\color{azulTitulo} ELECTRÓNICA --- Análisis de Imágenes}\\[4pt]
  {\large Informe de Análisis}\\[12pt]
  \rule{\linewidth}{1.2pt}\\[6pt]
  \rule{\linewidth}{0.6pt}
\end{center}

% ══ FICHA TÉCNICA ════════════════════════════════════════════════════════════
\vspace{4pt}
\begin{tcolorbox}[cajaTitulo, title=Datos del análisis]
\begin{tabular}{@{}llll@{}}
  \textbf{Fecha:}       & 2026-06-28  &
  \textbf{Modelo usado:} & gemini-2.5-flash \\[4pt]
  \textbf{Archivos:}    & \multicolumn{3}{l}{\texttt{imagen\_1.jpeg}, \texttt{imagen\_2.jpeg}} \\
\end{tabular}
\end{tcolorbox}

% ══ DESCRIPCIÓN GENERAL ══════════════════════════════════════════════════════
\section*{Descripción general}
Las imágenes muestran la pantalla de un ordenador con notificaciones de seguridad de Windows Defender. Se detectan varias amenazas clasificadas como 'HackTool:Python/Impacket', algunas de las cuales han sido puestas en cuarentena, mientras que otra permanece activa y ejecutándose en el dispositivo.

% ══ ELEMENTOS DETECTADOS ═════════════════════════════════════════════════════
\section*{Elementos detectados}
\begin{itemize}[leftmargin=1.5em, itemsep=2pt]
  \item \textbf{Interfaz de seguridad de Windows}: Ventana que muestra el historial de elementos recientes detectados por el antivirus.
  \item \textbf{Encabezado 'Todos los elementos recientes'}: Título de la sección actual en la interfaz de seguridad.
  \item \textbf{Botón 'Filtros'}: Opción para filtrar los elementos mostrados.
  \item \textbf{Alerta 'Amenaza encontrada: se requiere acción.'}: Notificación de alta prioridad con un icono de 'X' roja, indicando una amenaza que necesita intervención. Aparece dos veces, una con estado 'En cuarentena' y otra con estado 'Activo'.
  \item \textbf{Alerta 'Amenaza puesta en cuarentena'}: Notificación con un icono de escudo azul, indicando que una amenaza ha sido aislada. Se muestra con prioridad 'Alta'.
  \item \textbf{Detalles de la amenaza 'HackTool:Python/Impacket.Q'}: Información específica sobre una amenaza detectada, incluyendo su nombre, estado 'En cuarentena', fecha y hora (27/6/2026 10:32 a. m.), descripción del comportamiento potencialmente no deseado y el archivo afectado: 'C:\textbackslash\{\}Users\textbackslash\{\}E550\textbackslash\{\}Desktop\textbackslash\{\}CYBERSEGURIDAD\textbackslash\{\}venv\textbackslash\{\}Lib\textbackslash\{\}Scripts\textbackslash\{\}GetNPUsers.py'.
  \item \textbf{Detalles de la amenaza 'HackTool:Python/Impacket!AMTB'}: Información específica sobre otra amenaza detectada, incluyendo su nombre, estado 'Activo', indicando que no se corrigió y se está ejecutando. Fecha y hora (27/6/2026 10:32 a. m.), descripción del comportamiento potencialmente no deseado y los archivos afectados: 'C:\textbackslash\{\}Users\textbackslash\{\}E550\textbackslash\{\}Desktop\textbackslash\{\}CYBERSEGURIDAD\textbackslash\{\}venv\textbackslash\{\}Lib\textbackslash\{\}site-packages\textbackslash\{\}impacket\textbackslash\{\}spnego.py' y 'C:\textbackslash\{\}Users\textbackslash\{\}E550\textbackslash\{\}Desktop\textbackslash\{\}CYBERSEGURIDAD\textbackslash\{\}venv\textbackslash\{\}Lib\textbackslash\{\}site-packages\textbackslash\{\}impacket\textbackslash\{\}\_\_init\_\_.py'.
  \item \textbf{Enlace 'Más información'}: Opción para obtener detalles adicionales sobre la amenaza.
  \item \textbf{Botón 'Acciones'}: Botón desplegable para seleccionar acciones a tomar sobre la amenaza detectada.
  \item \textbf{Alerta 'Se bloqueó esta aplicación'}: Notificación con un icono de cuadrado azul, indicando que una aplicación fue bloqueada. Se muestra con prioridad 'Baja' y fecha 10/2/2026 10:50 p. m.
\end{itemize}

% ══ TEXTO EXTRAÍDO ═══════════════════════════════════════════════════════════
\section*{Texto extraído}
Seguridad de Windows
Todos los elementos recientes
Filtros
Amenaza encontrada: se requiere acción.
27/6/2026 10:32 a. m.
Alta
Amenaza puesta en cuarentena
27/6/2026 10:31 a. m.
Alta
Detectado: HackTool:Python/Impacket.Q
Estado: En cuarentena
Los archivos en cuarentena están en un área restringida donde no pueden dañar el dispositivo. Se quitarán automáticamente.
Fecha: 27/6/2026 10:32 a. m.
Detalles: Este programa tiene un comportamiento potencialmente no deseado.
Elementos afectados:
file: C:\textbackslash\{\}Users\textbackslash\{\}E550\textbackslash\{\}Desktop\textbackslash\{\}CYBERSEGURIDAD\textbackslash\{\}venv\textbackslash\{\}Lib\textbackslash\{\}Scripts\textbackslash\{\}GetNPUsers.py
Más información
Acciones
Baja
Se bloqueó esta aplicación
10/2/2026 10:50 p. m.
Detectado: HackTool:Python/Impacket!AMTB
Estado: Activo
No se corrigieron amenazas activas y se están ejecutando en el dispositivo.
Elementos afectados:
file: C:\textbackslash\{\}Users\textbackslash\{\}E550\textbackslash\{\}Desktop\textbackslash\{\}CYBERSEGURIDAD\textbackslash\{\}venv\textbackslash\{\}Lib\textbackslash\{\}site-packages\textbackslash\{\}impacket\textbackslash\{\}spnego.py
file: C:\textbackslash\{\}Users\textbackslash\{\}E550\textbackslash\{\}Desktop\textbackslash\{\}CYBERSEGURIDAD\textbackslash\{\}venv\textbackslash\{\}Lib\textbackslash\{\}site-packages\textbackslash\{\}impacket\textbackslash\{\}\_\_init\_\_.py

% ══ OBSERVACIONES ════════════════════════════════════════════════════════════
\section*{Observaciones}
Las imágenes muestran dos detecciones distintas de herramientas de hacking basadas en Python (Impacket). Una de ellas, 'HackTool:Python/Impacket.Q', ha sido exitosamente puesta en cuarentena, lo que indica que el sistema de seguridad actuó correctamente para neutralizarla. Sin embargo, la segunda detección, 'HackTool:Python/Impacket!AMTB', se encuentra en estado 'Activo' y el mensaje indica explícitamente que 'No se corrigieron amenazas activas y se están ejecutando en el dispositivo'. Esto representa un riesgo de seguridad crítico, ya que una herramienta de hacking está activa en el sistema. Ambas amenazas están asociadas a un entorno de desarrollo o proyecto llamado 'CYBERSEGURIDAD', lo que podría indicar que son herramientas legítimas utilizadas en un contexto de pruebas de seguridad, o que este entorno ha sido comprometido.

% ══ SUGERENCIAS ═════════════════════════════════════════════════════════════
\section*{Sugerencias}
La amenaza 'HackTool:Python/Impacket!AMTB' está activa y ejecutándose en el dispositivo, lo que representa un riesgo de seguridad inmediato. Se recomienda lo siguiente:
1.  **Tomar acción inmediata:** Hacer clic en el botón 'Acciones' junto a la amenaza activa y seleccionar la opción para 'Quitar' o 'Poner en cuarentena' la amenaza de forma manual.
2.  **Realizar un análisis completo del sistema:** Ejecutar un análisis completo con Windows Defender (o el antivirus principal) para asegurar que no haya otras amenazas persistentes o relacionadas.
3.  **Investigar el origen:** Dado que las amenazas se encuentran en la ruta 'C:\textbackslash\{\}Users\textbackslash\{\}E550\textbackslash\{\}Desktop\textbackslash\{\}CYBERSEGURIDAD\textbackslash\{\}venv\textbackslash\{\}Lib\textbackslash\{\}...', se debe investigar si estas herramientas son parte de un proyecto de ciberseguridad legítimo y si su ejecución es intencional. Si no es así, el entorno de desarrollo podría estar comprometido.
4.  **Desconectar de la red:** Si la amenaza persiste o no se puede eliminar, se recomienda desconectar el equipo de la red para evitar una posible propagación o exfiltración de datos mientras se resuelve el problema.

% ══ PIE DE INFORME ════════════════════════════════════════════════════════════
\vspace{12pt}
\begin{center}
  \footnotesize\color{gray}
  Informe generado automáticamente por el sistema de análisis con IA (gemini-2.5-flash) y soporte LaTeX. \\
  Generado el 2026-06-28 \textbullet\ ELECTRÓNICA --- Sistema de Análisis de Imágenes.
\end{center}

\end{document}
